% FILE: 03_architecture_and_primitives.tex
% PROJECT: otel-weaver-exp-001-custom-pi-registry
% THESIS ROLE: otel-semconv-evidence-surface-experiment

\section{Architecture and Primitives}
\label{sec:otel-weaver-exp-001-custom-pi-registry-architecture}

\subsection{Core Objects}
\label{sec:otel-weaver-exp-001-custom-pi-registry-architecture-objects}

The \texttt{process.pi.activity} span group is the central schema object. It organizes eight
named attributes into two cardinality classes:

\textbf{Required attributes:}
\begin{itemize}
  \item \texttt{process.pi.instance\_id} --- the case identifier; maps to the OCEL case ID that
        groups all events belonging to one process instance.
  \item \texttt{process.pi.activity.name} --- the Petri net transition label; must name a
        declared transition in the process model, not a free-form application string.
  \item \texttt{process.pi.lifecycle} --- the lifecycle phase; constrained to the enumeration
        \{\texttt{schedule}, \texttt{start}, \texttt{suspend}, \texttt{resume},
        \texttt{complete}, \texttt{abort}\}; maps to OCEL lifecycle transitions.
  \item \texttt{process.pi.witness.hash} --- a BLAKE3 64-character hexadecimal seal over the
        span's process-relevant payload; provides independent verifiability.
  \item \texttt{process.pi.witness.id} --- the identifier of the witness node that sealed the
        span; enables tracing of witness provenance across distributed executions.
\end{itemize}

\textbf{Recommended attributes:}
\begin{itemize}
  \item \texttt{process.pi.token.state\_before} --- JSON-serialized Petri net marking before the
        transition fires; enables pre-condition verification in replay.
  \item \texttt{process.pi.token.state\_after} --- JSON-serialized Petri net marking after the
        transition fires; enables post-condition verification in replay.
  \item \texttt{process.pi.activity.type} --- an architectural category tag; recommended, not
        required; its optional status creates a known open gap (see Section~\ref{sec:otel-weaver-exp-001-custom-pi-registry-questions-gaps}).
\end{itemize}

The downstream Rust structs \texttt{BridgeRx} (EXP-002, \texttt{bridge.rs}) and the
\texttt{validate\_and\_project} function (EXP-003, \texttt{validator.rs}) consume these attributes
as their input contract. The zero-sized witness type
\texttt{OtelSpanToOcelEventProjection} seals the translation from OTel span to OCEL event.

\subsection{Admissible Transitions}
\label{sec:otel-weaver-exp-001-custom-pi-registry-architecture-transitions}

A span transitions from \emph{raw telemetry} to \emph{admissible feedstock} by satisfying all
five required attributes. The transition is gated by the \texttt{validate\_and\_project} function
in EXP-003, which produces either an \texttt{Admission<T,W>} value (feedstock accepted, typed by
the witness \texttt{W}) or a \texttt{Refusal} value carrying a \texttt{RefusalCode} from the
enum: \{\texttt{MissingInstanceId}, \texttt{MissingActivityName}, \texttt{MissingWitnessId},
\texttt{InvalidWitnessSignature}\}. The gate is a type-law gateway: refusal is expressed as a
typed value in the Rust type system, not as a runtime exception or a logged warning.

Once admitted, feedstock passes through \texttt{BridgeRx} for schema-version translation. The
bridge translates schema-version diffs without producing residual state: a span arriving from a
registry version prior to the current pin is bridged to the current attribute set without
contaminating the conformance court with legacy field interpretations. This property --- no
residual from bridging --- is the \emph{BridgeRx} invariant stated in EXP-002's \texttt{bridge.rs}.

\subsection{Boundaries}
\label{sec:otel-weaver-exp-001-custom-pi-registry-architecture-boundaries}

EXP-001 is bounded on the upstream side by the OTel instrumentation layer: the experiment
specifies the schema that instrumented services must emit but does not specify how instrumentation
is implemented. The experiment is bounded on the downstream side by the \texttt{BridgeRx}
layer in EXP-002: the registry defines the attribute set; EXP-002 handles version translation;
EXP-003 handles admission gating. EXP-001 does not perform conformance replay; that responsibility
belongs to \texttt{wasm4pm} and \texttt{wasm4pm-compat} after feedstock has been admitted and
projected to OCEL.

The doctrine boundary is stated in \texttt{otel-weaver/doctrine/otel-weaver-is-feedstock.md}:
OTel Weaver is a feedstock-layer tool, not a process-intelligence tool. It validates that spans
satisfy a schema; it does not determine whether those spans conform to a process model. The
process-intelligence claim begins after the feedstock boundary, not at it.

\subsection{Failure Modes Refused}
\label{sec:otel-weaver-exp-001-custom-pi-registry-architecture-failures}

The architecture explicitly refuses the following failure modes:

\begin{enumerate}
  \item \textbf{Schema-drift-as-process-drift conflation.} A change to \texttt{process\_pi.yaml}
        does not imply that the process has changed. Registry diff tracking and process conformance
        tracking are independent concerns, as stated in
        \texttt{doctrine/registry-diff-is-not-process-drift.md}.
  \item \textbf{Telemetry-as-evidence conflation.} A span that satisfies no required attributes
        is not process evidence. It is refused at the \texttt{validate\_and\_project} gate with
        a typed \texttt{RefusalCode}; it does not silently pass into the conformance court.
  \item \textbf{Aspirational receipt claims.} The receipt chain is declared in the manifest but
        is marked PARTIAL until the file at \texttt{receipts/otel-weaver-source-receipt.json} is
        materialized on disk. The architecture refuses to treat a declared receipt as an
        equivalent of a sealed receipt.
  \item \textbf{Self-certified ALIVE status without independent replay.} The ALIVE checkpoint
        was issued by the ggen manufacturing authority. The architecture treats this as a
        necessary but not sufficient condition for full ALIVE status; an independent pm4py-based
        fitness/precision check against a live event log is required to close the PARTIAL gap.
\end{enumerate}
